\section{Universal covers and the uniform simplex metric}
\label{sec:closed-covers}

We use the continuous macroscopic dimension of a metric space $X$:
$\dim_{\mathrm{mc}}X\le d$ means that a continuous map from $X$ to a
simplicial complex of dimension at most $d$, with its uniform simplex
metric, has uniformly bounded entire fibers.  Here each simplex is a
standard Euclidean simplex with the same edge length.  This is the
continuous convention in \citet[Definitions~1.1--1.2]{dranishnikov2014}.
We first convert the map already constructed to this target metric.
We then supply the two lower dimensions using the classical geodesic
index argument and the width theorem of Liokumovich and Wang.

\begin{lemma}[Changing the simplex metric without changing fibers]
\label{cover:uniform}
Let $k\ge0$ be an integer.  Suppose continuous functions $t_i:Y\to[0,1]$ have locally finite
supports, at least one equals $1$ at each point, and at most $k$ are
strictly between $0$ and $1$.  Their coordinate map admits an injective
conversion to a simplicial complex of dimension at most $k$ with its
uniform simplex metric.  The converted map is continuous, and two of
its values agree if and only if all the original coordinates agree.
\end{lemma}
\begin{proof}
Use the cubes $Q(S,A)$ and staircase simplices of
\eqref{topology:staircase}, where $S,A$ are disjoint finite sets,
$S\ne\varnothing$, and $|A|\le k$.  Realize the abstract simplices
as standard simplices, and send each coordinate simplex to its new
realization with the same barycentric coordinates.  The rules agree
on common faces.  Equality of two images identifies their minimal
simplex and its positive barycentric coordinates; the affine
combination of the corresponding finite indicator vectors recovers
the same original coordinate vector.  Thus the conversion is injective.

For $k\ge1$, give each simplex edge length $1$ and define $d_K$
as the infimum of the total Euclidean lengths of finite piecewise
straight paths in the simplices.  This infimum is finite: vertices
$\mathbf1_S$ and $\mathbf1_T$ are joined by adding the coordinates
in $T\setminus S$ and then removing those in $S\setminus T$, always
retaining a nonempty set of coordinates equal to $1$.  Every point
is joined to a vertex in its simplex.  For a vertex $v$, its
barycentric coordinate $b_v$, set to zero on simplices not containing
$v$, agrees on faces and is $\sqrt2$-Lipschitz on each unit-edge
simplex.  Indeed such a simplex is realized by the vectors
$e_v/\sqrt2$ in a Euclidean coordinate space.  Hence
\[
 |b_v(x)-b_v(y)|\le\sqrt2\,d_K(x,y).
\]
The barycentric coordinates distinguish points, so $d_K$ is a genuine
finite metric.  For $k=0$ use the discrete vertex metric.

Near a point of $Y$, only a finite set $F$ of coordinates can be
nonzero.  The coordinate map there lands continuously in the finite
cubical subcomplex with all other coordinates zero.  The conversion
is continuous on each of its closed simplices and hence on the finite
complex.  Inclusion into the full uniform realization is continuous
on a simplex because distance there is at most the length of its
straight segment.  Finite pasting therefore proves continuity on
the entire finite subcomplex, and thus near every point of $Y$.
The same local finite argument applies for $k=0$.  No equality between two topologies on an arbitrary
infinite complex is required.
\end{proof}

\begin{lemma}[Level components of a proper Morse function]
\label{cover:reeb}
Let $u:N\to\R$ be a proper Morse function on a connected smooth
manifold without boundary.  There is a locally finite simplicial graph $K$ and a
continuous map $q:N\to K$, with $K$ carrying its uniform simplex
metric, whose point fibers are precisely the connected components
of the levels of $u$.
\end{lemma}
\begin{proof}
Choose an auxiliary smooth Riemannian metric on $N$ to define gradients.
The critical set is closed and its points are isolated.  Properness
therefore makes the number of critical points in any compact slab
$u^{-1}([a,b])$ finite.  Thus the critical values form a locally finite
subset of $\R$, even if several critical points have the same value.
Every level is compact and locally path connected: regular charts
give a hypersurface, while the Morse normal form at a critical point
gives the zero set of a nondegenerate quadratic form, which is a cone
or an isolated point.  A compact locally path connected space has finitely
many components, since its components are open and form an open cover.

Let $S$ consist of the critical values and the integers.  It is
locally finite.  At each $a\in S$ put a vertex for every component
of $u^{-1}(a)$.  Between consecutive values $a<b$ of $S$, the flow of
$\nabla u/|\nabla u|^2$ identifies $u^{-1}((a,b))$ with a regular
level times $(a,b)$.  The flow exists for each compact subinterval
because its trajectory stays in a compact slab.  Each of the finitely
many connected product components defines an open edge, parametrized
by its $u$-value.  The image of a proper map to $\R$ is closed;
any finite endpoint of $u(N)$ is an attained extremum and hence a
critical value.  Thus no further finite endpoints are missing.

An edge component $E$ has a unique attachment at each end.  At its
lower end, the nonempty compact connected sets
\[
 \overline{E\cap u^{-1}((a,a+\epsilon])},
 \qquad 0<\epsilon<b-a,
\]
are nested as $\epsilon\downarrow0$.  Their intersection equals
$\overline E\cap u^{-1}(a)$: a point of this latter set is approached
by points of $E$ whose heights tend to $a$, and therefore belongs
to every displayed closure.  The intersection is nonempty,
connected and contained in $u^{-1}(a)$, hence in exactly one of that
level's components.  Attach the edge to its vertex.  The same argument
applies at $b$.  There are finitely many edges in either adjacent slab,
so every vertex has finite valence.  The resulting graph is countable
and locally finite.  Its points correspond exactly to the connected
components of all levels, using full level components as vertices
when critical values coincide.

Define $q$ by this correspondence.  On an open edge it is continuous
by the product description.  At a vertex corresponding to a compact
level component $C\subset u^{-1}(a)$, an adjacent edge not attached
to $C$ has no sequence approaching $C$: such a sequence would put
a point of $C$ into its endpoint intersection, contrary to its unique
attachment.  The finitely many other components of $u^{-1}(a)$ are
disjoint compact sets.  A neighborhood of $C$ can therefore exclude
them and the finitely many nonincident edges.  Requiring also that
$|u-a|$ be sufficiently small places its remaining points in any
prescribed short tails of the finitely many incident edges.  This
proves continuity at the vertex and shows that its whole inverse
image is $C$. The graph is connected, since $q$ is a continuous
surjection and $N$ is connected.

Subdivide every edge once with a different midpoint vertex for each
edge.  This removes parallel edges from the abstract realization;
loops do not occur because the endpoint heights differ.  The result
is a simplicial graph.  Give each new edge the same length.  Finite
valence implies that its ordinary graph topology equals the topology
of this uniform simplex metric: vertex neighborhoods contain short
segments on each of finitely many incident edges.  The subdivision
and change of edge lengths preserve points and fibers.  Thus $q$
has the asserted continuity and exact fibers.
\end{proof}

\begin{proof}[Proof of \cref{cover:main}]
Put $s=\min_M\Scal_h>0$ and $X=\widetilde M$.  The lifted metric
is complete.  If $X$ is compact, use a constant map to a vertex,
whose only nonempty fiber has finite diameter.

For $n\ge4$, the complete metric $s\widetilde h$ has scalar
curvature at least $1$.  Apply \Cref{thm:main} using the coordinates
of \eqref{topology:coordinates}, then \Cref{cover:uniform}.  Local
finiteness of their supports was proved in
\Cref{sec:reconstruction}, and the conversion preserves their
entire fibers.  Scaling back gives the bound $C_n/\sqrt{s}$ in
the original metric.

For $n=2$, the Gaussian curvature of the cover is at least $s/2$.
If a unit-speed minimizing geodesic had length $L>\pi\sqrt{2/s}$,
take a parallel unit normal $E$ along it and the endpoint-zero
variation field $V(t)=\sin(\pi t/L)E(t)$.  Its index form obeys
\[
 I(V,V)=\int_0^L\bigl(|\nabla_tV|^2-K|V|^2\bigr)\,dt
 \le \frac{\pi^2}{2L}-\frac{sL}{4}<0,
\]
contradicting minimality.  Completeness supplies minimizing geodesics,
so the cover has diameter at most $\pi\sqrt{2/s}$ and is compact
by Hopf--Rinow.  The constant map therefore proves this case.

Finally let $n=3$ and $X$ be noncompact.  Apply
\citet[Theorem~1.1, arXiv version~1]{liokumovichWang2023} to the
complete metric $(s/6)\widetilde h$, whose scalar curvature is at
least $6$.  It gives a proper Morse function such that every
connected component of every level, including critical levels, has
ambient diameter less than $18\pi$ in that metric.  Apply
\Cref{cover:reeb}.  Each whole fiber of its graph map is one such
level component, so in the original lifted metric it has diameter
less than $18\pi\sqrt{6/s}$.  This is a continuous map to a
one-dimensional uniform simplicial complex, as required.
\end{proof}

The low-dimensional conclusions use their stated classical and public
inputs.  The quantitative theorem in dimensions at least four has
the separate proof given above.  All map statements here are continuous;
changing the simplex metric does not supply a Lipschitz bound.
